% Threats to validity, arranged as Roel Schuurkes (2026Oct/0161, 0163) and
% Guigui Wang (2026Oct/0148) left to the editors. Their texts are quoted
% whole between their own scissor marks; tests/test_sections_tex.py checks
% the quotes word for word against quoted/*.txt. Per 2026Oct/0147 this
% arrangement goes to Roel, Rouiada, Lars and Guigui before the draft.
\providecommand{\aevmsg}[2]{\href{https://lists.w3.org/Archives/Public/public-agent-conformance/2026#1/#2.html}{#1/#2}}

\section{Threats to validity}
\label{sec:w3c-threats}

Roel Schuurkes wrote that his limitations text and Guigui Wang's reader-side text ``would sit beside'' each other and that ``the outline can decide how the two divide'' (\aevmsg{Oct}{0161}); his text ends ``How this text relates to Guigui's reader-side text in 0148 is for the editors'' (\aevmsg{Oct}{0163}), and Guigui's begins ``Section number and placement are for the editors'' (\aevmsg{Oct}{0148}). This section proposes that arrangement. Each part names its author; the two posted texts are quoted whole, and the dated editorial update and the construct, external and conclusion parts add text.

\subsection{Internal validity (Roel Schuurkes, Oct/0163)}

\begin{quote}
What the table establishes is narrow. Where a row reports a rerun, it is a reproduction in O8's sense: the producer's implementation, run by another operator with the expected outputs visible. It shows that published bytes are what a pinned commit produces. It does not show that a checker, an expected rule or a reading of the text is right. The readers written from the text alone (R8, R11) are the exceptions, and each is one builder's run whose ``from the text alone'' rests on the exposure that builder recorded. In R4, two readers were written from the published rules by their own authors but ran through the producer's glue with the answers visible, so they come close to an independent record without being one.

The units differ by row: checks, members, mutants, synthetic cases. Counts from different rows should not be added or turned into a common rate.

As 0148 says, every quantitative claim in the paper must name its corpus pin. In the table, R8 and R10 are at v0.17.4, R9 compares v0.17.4 with v0.18.0 and v0.18.1, R11 and R12 are at v0.18.1, and R13 carries its own commit pins.

A comparison and the run it compares are separate evidence. The 92-of-96 figure in R8 is a reviewer's comparison that neither rerun repeats, and the causes in R11 were written by the session that set up the build, which had seen the expected outcomes.

Rows that rest on text that v0.1 does not hold can be reported as proposed, not as conformance to v0.1: the 5.5 fixture records (R2, R3), the jcs-sha256 members (R12) and W3C-R-033 (R13).

R14 lists only reader verdicts posted on the list. None has been posted for R3 or R6, so the table says nothing about whether today's readers can check those two reports.

The run outputs of my own reruns are kept locally; only what the messages state is on the record, as 0061 and 0064 say for R3 and R2.
\end{quote}

\paragraph{Editorial update, 3 November 2026.} The preceding quotation preserves the limitations posted on 23 October (\aevmsg{Oct}{0163}). Since then, the producers have posted converted reports for R3 and R6 (\aevmsg{Nov}{0000}, \aevmsg{Nov}{0006}); the committed matrix records R3-v01 and R6-v01 separately from their original inputs. Acceptance of a conversion is not retrospective acceptance of the original bytes. The quotation's statement about the R8 comparison describes the reruns available at that time. The later R8-compare record (\aevmsg{Oct}{0188}, \aevmsg{Nov}{0002}) remains a comparison, with its own selected populations and exposure history. These updates do not turn a reader's acceptance into evidence that the producer's checks are correct, or claim a new execution for this editorial revision.

\subsection{Construct validity (Guigui Wang, Oct/0148)}

\begin{quote}
Three readers, three evidentiary classes. The runs the paper relies on involve three kinds of reader, and the limits section should keep their evidentiary weight apart: (1) a reader written from the v0.1 text alone, importing none of the reference implementation, with trust starting at the release cosign.pub; (2) a reader built from the text with the reference source in view; (3) the reference reader itself. Kind (2) agreeing with kind (3) is not independence. In one run, a reader written with the reference source in view returned 235/235 agreement with the reference reader on the v0.17.4 report corpus, including the W3C-R-016 rejection of member v04b7811e691bd2c9; v0.18.0 later withdrew that rejection with forty-one others. That 235/235 agreement is therefore reproduction and, at the same time, evidence against independence: the reader reproduced the reference's mistake.

What a text-alone reader can decide. On questions the text determines, an independent reader recomputes from released bytes and reaches the answer without reference code: cryptographic signature verification (22/25 receipt signatures verified, with the remaining three failing exactly as the corpus marks them), semantic verdicts on a suite whose expected outcomes are fixed (25/25 on the receipt-signature suite), the signed CORPUS-DIGESTS chain over the listed corpora, and serialization questions the governing canonical form settles. On integer minus-zero, the Python reference reader wrote 0 and the Go reader -0 at v0.17.4 (0038); an independent reader from the text returned RFC 8785's answer, 0 (0046), and no verdict depended on that member. Distribution gaps are visible to such a reader as well: one listed corpus ships in git only and is absent from the published wheel, so a reader trusting released bytes and the signed digest list is presented with a manifest it cannot resolve.

Where the text leaves the answer open. On the forty-three members rejected under W3C-R-016, a text-alone reader's step reports held both readings of the section 2 rule (S-15) - moved value as declared versus moved recomputed - and selected neither, answering cannot-tell (0039; 0098/0101/0102). The reference reader read the declared value (0054); v0.18.0 moved forty-two of those members to accept (0079) and re-cut the forty-third, a w3c-f-16 member, with its accept twin (0103). The section-10 selection changed from 83 reject / 82 accept to 41 reject / 124 accept while the v0.1 text itself did not change. The reader-side conclusion is twofold. First, cannot-tell is a finding: it locates an ambiguity in the text rather than silently accepting or silently rejecting. Second, every quantitative claim in the paper must name its corpus pin, because the same reader over adjacent corpus releases gives materially different counts; unpinned reject/accept percentages are not comparable across releases.

Reproduction hierarchy. The frozen re-runs fall into distinct classes, and the paper should label each row. Bit-for-bit operator reproduction of the text-alone reader over all 235 v0.17.4 members (corpus commit 99b47c56...) matched verdicts.tsv (de987220..., 169 lines) and not-selected.tsv (cd6962f3..., 66 lines) on two independent environments, Linux arm64 and Linux x86\_64 with CPython 3.13.13 and PYTHONHASHSEED=0. A cross-environment re-run of a pinned implementation packet at d0cbe773 (CPython 3.13.12, cryptography 46.0.5, fresh checkout, PYTHONHASHSEED=0) reproduced its 18 tests and 21 selected outcomes (7 pass / 13 fail / 1 inconclusive) with both file digests matching; this is category (a) under O8 - same implementation, independent operator and environment - not an independent implementation. No rerun that shares the scorer can establish that the expected verdicts are right, because the expected answers were in view. Establishing the readings themselves requires readers built without reference code, and the corpus-mutation history between v0.17.4 and v0.18.0 is the strongest evidence the group has for why that distinction matters.
\end{quote}

\paragraph{What a conforms verdict establishes (the appendix's author).} A conforms verdict establishes that the report holds against the rows the named reader release checks, at that release; it does not establish that the producer's checks measure what their names say. The verification section's generated matrix gives the stored negative-control outcomes and denominators per reader. Its positive controls and targets marked not applicable are separate populations. The readers written from the text can answer cannot-tell where the text leaves a reading open (\aevmsg{Oct}{0039}); their stored answers are findings for their authors, not automatically faults. The matrix compares readers on pinned report bytes and remains distinct from an independent implementation of the producer's checks.

\subsection{External validity (the appendix's author)}

Every producer in the table chose to take part, and most helped write or adopt the format they report in: the reference emitter (R1) and the reference readers are the appendix author's, and the other producers posted their reports to the group that wrote v0.1. Two rows rest on generated data in their owners' own words: R5 counts ``one generated draw against one threshold'' over 5,000 mock draws per rail, six thresholds per draw plus six exclusions (\aevmsg{Oct}{0162}, \aevmsg{Nov}{0020}), and R13 runs five synthetic cases (\aevmsg{Oct}{0162}). All of it comes from one community group between September and November 2026. A producer that did not help write the format is the test it lacks: Evgenii Arsentev offered ``a second producer (reports from my own agent runs), so the format meets data it wasn't built around'' (\aevmsg{Oct}{0005}), and no such report has been posted yet.

\subsection{Conclusion validity (rule: Roel Schuurkes, Oct/0163)}

The rule is Roel Schuurkes's: ``Counts from different rows should not be added or turned into a common rate'' (\aevmsg{Oct}{0163}). The verification table gives each row its unit and denominator so that the rule can be checked row by row; no figure in this paper sums or averages across rows.
